\chapter{Technologie i narzędzia}

Istotny był wybór odpowiednich technologii i narzędzi. Ważne było aby system był wydajny i jego poszczególne komponenty łatwe do wdrożenia dzięki odpowiednim udokumentowaniu. Celem tego rozdziału jest przedstawienie użytych w projekcie technologii oraz uzasadnienie ich wyboru.

\section{Wybrane języki programowania}

Do realizacji projektu zdecydowano się na język programowania Python, który dzięki swojej prostocie, szybkości implementacji oraz sporej liczbie bibliotek, był najlepszym wyborem z pośród innych języków. Dodatkowym atutem jest jego popularność pośród istniejących już rozwiązań uczenia maszynowego. Stanowiło to dodatkową inspirację podczas projektowania systemu.

\section{Biblioteki}

W projekcie zostały wykorzystane różne biblioteki, które pozwoliły przyspieszyć czas implementacji, zmniejszając potrzebę pisania kodu od podstawy oraz utrzymania wysokiej jakości rozwiązań problemów programistycznych.

- json
- random
- re
- requests
- subprocess

\section{Wybrane modele językowe}

Po szeregu różnych testów zdecydowano się na wykorzystanie modeli

\section{Narzędzia}

\subsubsection{Git}

Git to narzędzie służące do zarządzania wersjami projektu, umożliwiające monitorowanie wprowadzanych zmian w kodzie oraz równoczesną pracę wielu osób nad tym samym projektem.

W projekcie pełnił funkcję podstawowego systemu kontroli wersji. Każdy członek zespołu pracował na oddzielnej gałęzi (ang. branch), a po zakończeniu prac wprowadzone zmiany były łączone z główną gałęzią projektu (ang. main branch). Takie podejście ułatwiało współpracę, ograniczało ryzyko występowania konfliktów oraz pozwalało na zachowanie uporządkowanej, liniowej historii zmian.

\subsubsection{GitHub}

GitHub to platforma wykorzystująca system Git, która pozwala na przechowywanie oraz zarządzanie repozytoriami znajdującymi się na zdalnych serwerach. Oprócz podstawowych funkcji oferuje również możliwość wzajemnego przeglądania i weryfikowania kodu przez członków zespołu, a także automatycznego uruchamiania procesów sprawdzających poprawność projektu, takich jak wykonywanie testów czy weryfikacja kompilacji kodu w ramach CI (ang. Continuous Integration). Funkcja przeglądu kodu pozwalała zespołowi na bieżąco kontrolować jakość wprowadzanych zmian, wykrywać potencjalne błędy oraz utrzymywać kod w odpowiednim stanie.

\subsubsection{Ollama API}

Ollama to oprogramowanie umożliwiające lokalne uruchamianie oraz zarządzanie wybranymi modelami językowymi typu LLM (ang. Large Language Model).

W projekcie zostało wykorzystane do obsługi modeli działających bezpośrednio na lokalnej infrastrukturze, co pozwalało na korzystanie z ich możliwości bez konieczności przesyłania danych do zewnętrznych usług. Komunikacja pomiędzy aplikacją a uruchomionymi modelami odbywała się za pośrednictwem interfejsu API, umożliwiając przesyłanie zapytań oraz odbieranie wygenerowanych odpowiedzi.

\section{Wymagania sprzętowe}